% ============================================================================
\section{Mechanism: Visible Losses Control Visible Coordinates}
\label{sec:mechanism}

The failure mode is a mismatch between what the recurrent state carries and what the loss can see.
At loop $k$, a hidden state $H_k$ is both a prediction interface and the input to future computation:
\begin{equation}
    H_{k+1}=F_\theta(H_k),\qquad z_k=W_{\mathrm{out}}r(H_k),\qquad
    \mathcal{L}_k=\mathrm{CE}(z_k,y).
\end{equation}
Here $F_\theta$ is the recurrent transition, $r$ is the readout, $W_{\mathrm{out}}$ maps readout states to logits $z_k$, $y$ is the next-token target, and $\mathcal{L}_k$ is the loop-$k$ cross-entropy loss.
A degree of freedom is \emph{immediately visible} if the local readout loss has a derivative in that direction, and \emph{active} if the transition updates it and carries it forward.
Per-loop CE makes each loop predictive, but it does not make every active degree of freedom visible.

We study the scale direction.
Let $H=su$, where $s=\|H\|_{\rms}$, $\|u\|_{\rms}=1$, and $\|v\|_{\rms}^2=d^{-1}\|v\|_2^2$ for hidden-vector dimension $d$.
A raw readout preserves scale, $r_{\mathrm{raw}}(su)=su$.
A normalized readout removes it, $r_{\mathrm{norm}}(su)=u$ up to the usual normalization epsilon.
Figure~\ref{fig:mechanism} sketches the two paths the same state takes.

\begin{figure}[t]
\centering
\includegraphics[width=\linewidth]{figures/visibility_activity_mechanism.pdf}
\caption{\textbf{Visibility--activity mismatch.}
The same hidden state $H_k=s_ku_k$ serves as both a prediction interface and the input to the next recurrent step.
Along the readout path, output normalization removes scale, so the immediate CE loss has approximately zero radial gradient.
Along the recurrent path, the pre-norm residual update still carries the skip state $H_k$, so scale remains active and can drift through the recurrence.
Here $\mathrm{Norm}$ denotes the pre-sublayer normalizer and $B$ denotes the residual branch applied to the normalized state.}
\label{fig:mechanism}
\end{figure}

\begin{lemma}[Visibility: scale-invariant readouts remove the immediate radial CE signal]
\label{lemma:radial}
If $r(\alpha H)=r(H)$ for all $\alpha>0$ and $\mathcal{L}(H)=\mathrm{CE}(W_{\mathrm{out}}r(H),y)$ is differentiable at $H\neq0$, then
\begin{equation}
    \langle \nabla_H \mathcal{L}(H),H\rangle=0.
\end{equation}
\end{lemma}
\noindent
The proof is the chain rule applied to $\mathcal{L}(e^tH)$, which is constant in $t$.
With practical LayerNorm/RMSNorm epsilons the identity is approximate, but the radial derivative is small at large hidden norm.

\begin{lemma}[Activity: pre-norm residual loops carry scale]
\label{lemma:expansion}
Let $H=su$ with $\|u\|_{\rms}=1$.
For the epsilon-free pre-norm residual update, with $\mathrm{Norm}$ the pre-sublayer normalizer and $B$ the residual branch,
\[
    F(H)=H+B(\mathrm{Norm}(H)),
    \qquad
    b(u)=B(\mathrm{Norm}(u)),
\]
we have $F(su)=su+b(u)$.
Decompose the residual update into radial and angular parts:
\[
    b(u)=a_{\rad}(u)u+b_\perp(u),
    \qquad
    a_{\rad}(u)=\frac{\langle u,b(u)\rangle}{d},
    \qquad
    \langle u,b_\perp(u)\rangle=0.
\]
If $\|b(u)\|_{\rms}=O(1)$ as $s\to\infty$, then
\begin{align*}
    \text{\emph{scale update:}}\qquad
    \|F(su)\|_{\rms}
    &= s+a_{\rad}(u)+O(s^{-1}), \\
    \text{\emph{direction update:}}\qquad
    \frac{F(su)}{\|F(su)\|_{\rms}}
    &= u+\frac{b_\perp(u)}{s}+O(s^{-2}).
\end{align*}
\end{lemma}
\noindent
Thus normalized readouts remove the direct radial CE signal, while pre-norm recurrence can still update scale at leading order.
The lemmas identify a local mechanism, not an unconditional theorem of explosion.
The sign and persistence of $a_{\rad}(u_k)$ are learned empirical properties.
If the learned block has persistent positive radial velocity, however, the local CE path through a normalized readout has no direct way to oppose it.
Combining the two statements gives the local scale-control mismatch
\begin{equation}
    \langle \nabla_{H_k}\mathcal{L}_k,H_k\rangle=O(\epsilon/s_k^2),
    \qquad
    s_{k+1}-s_k=a_{\rad}(u_k)+O(s_k^{-1}),
\end{equation}
where $\epsilon$ is the normalizer epsilon; the first term is exact in the epsilon-free idealization and approximate with practical normalizers.
Proof details are in Appendix~\ref{app:proofs}.

\paragraph{Tokenwise normalization and transformer blocks.}
The scalar decomposition keeps notation readable.
In the actual transformer, RMSNorm is applied over channels separately at each token, so the scale-blind direction is one radial degree of freedom per token vector.
The same expansion also applies to a multi-sublayer pre-norm decoder block at leading order: after the first residual addition, later normalizers see $su+O(1)$, so their normalized input differs from $\mathrm{Norm}(u)$ only by $O(s^{-1})$.
Appendix~\ref{app:per-token-norms} reports token-level norm statistics.

\paragraph{Indirect future-loop signals.}
Future losses are not zero-signal.
A later loss can backpropagate through the recurrence and therefore depend indirectly on the earlier scale.
The issue is attenuation: for a fixed future horizon $t$, changing $s_k$ mainly changes the size of later angular steps, $u_{j+1}-u_j=b_\perp(u_j)/s_j+O(s_j^{-2})$.
Under smoothness and bounded-update assumptions on a stable trajectory, this yields
\begin{equation}
    \biggl|\frac{\partial \mathcal{L}_{k+t}}{\partial \log s_k}\biggr|
    \lesssim \frac{C t}{s_k}
\end{equation}
for a constant $C$ depending on those bounds.
We offer this as an explanatory estimate rather than a global-training theorem: it is the fixed-depth reason future CE can be a weak scale-control path at large norm.

\paragraph{Design rule.}
Dense CE trains visible prediction interfaces.
Recurrent scale control requires either making scale visible to at least one loss, for example with raw readouts or an explicit norm penalty, or removing/resetting scale inside the recurrent path, as inter-loop normalization does.
Table~\ref{tab:interventions-main} organizes the interventions used in this paper by which of these roles they play.

\begin{table}[t]
\centering
\small
\caption{\textbf{Scale control and exit training are separate design requirements.} The interventions separate whether CE sees scale, whether scale is controlled by an auxiliary loss, or whether scale is removed from the recurrent path.}
\label{tab:interventions-main}
\begin{tabular}{p{0.30\linewidth}p{0.60\linewidth}}
\toprule
\textbf{Intervention} & \textbf{Role in the framework} \\
\midrule
\multicolumn{2}{l}{\emph{Scale-hidden CE}} \\
Normalized readout & Hides hidden-state scale from the immediate CE loss. \\
Per-loop CE through normalized readouts & Trains intermediate prediction exits, but does not by itself make scale visible. \\
\addlinespace
\multicolumn{2}{l}{\emph{Scale-visible CE}} \\
Raw readout & Makes scale visible to CE by breaking scale invariance. \\
\addlinespace
\multicolumn{2}{l}{\emph{Explicit scale penalty}} \\
Norm penalty & Adds an auxiliary scale-visible objective while keeping the normalized readout. \\
\addlinespace
\multicolumn{2}{l}{\emph{Hybrid readout}} \\
Final-only norm & Exposes intermediate loops through raw readouts but keeps a normalized final interface. \\
\addlinespace
\multicolumn{2}{l}{\emph{Scale-removing recurrence}} \\
Inter-loop normalization & Removes scale from the recurrent input rather than supervising it. \\
\bottomrule
\end{tabular}
\end{table}

Raw readout by itself is not the causal proof.
It is a useful scale-visible intervention, but it also changes logit scale, calibration, output-weight usage, and gradient magnitudes.
The scale-visibility interpretation comes from convergence across controls: raw readouts, explicit norm penalties, final-only normalization, radial-gradient diagnostics, and the recurrent-scale clamp all point to the same active-but-hidden scale direction.
